% Standalone appendix document - compiles on its own:
%   latexmk -pdf -outdir=build 04-meeting-minutes-2026-10-06.tex
\documentclass[a4paper,11pt]{article}
\usepackage[danish]{babel}
\usepackage{../../appendix-style}

\title{13.4 Mødereferat -- 06-10}
\author{SW4PRJ4 -- Group 3}
\date{06-10-2026}

\begin{document}
\maketitle

\section*{Mødeinformation}

\textbf{Dato:} 06-10-2026 (tirsdag)

\textbf{Tid:} 13:00

\textbf{Sted:} 5123-313

\textbf{Referent:} Frederik

\section*{Deltagere}

\begin{itemize}
    \item Alexander H
    \item Alexander D
    \item Frederik
    \item Ida D
    \item Ida F
    \item Michel (vejleder)
\end{itemize}

\textbf{Fraværende:}

\begin{itemize}
    \item Malthe
    \item Marie
    \item Oskar
\end{itemize}

\section*{Dagsorden}

\begin{enumerate}
    \item Status siden sidste vejledermøde (15/9)
    \item Spørgsmål til vejleder
    \begin{enumerate}
        \item Diagrammer
        \item Kilder til kursusmateriale
    \end{enumerate}
    \item Næste vejledermøde
    \item Eventuelt
\end{enumerate}

\section*{Noter}
\subsection*{1. Status siden sidste vejledermøde}

\begin{itemize}
    \item \textbf{Deployment af kunde-appen:} Vi skal beskrive, at appen udgives gennem Expo og ikke direkte til Google Play.
    \item \textbf{Auth-flowet:} Hele flowet fra mobil til Keycloak og videre til backend skal beskrives ordentligt. Kunde-appen og admin-panelet kan have forskellige flows, og i så fald skal begge dokumenteres.
    \item Michel vil gerne have et flowdiagram over auth-flowet med teknisk tekst på pilene, så man kan se, hvad der faktisk bliver sendt mellem parterne.
    \item Credentials sendes kun mellem klienterne og Keycloak, aldrig gennem backend. Backend håndterer kun tokens fra klienterne og læser bl.a. claims i JWT-tokenet.
    \item Vi skal finde ud af, hvilke claims de forskellige brugere skal have. Et claim som ``onboarding gennemført'' kan spare et opslag i databasen.
    \item Roller og rettigheder ligger typisk også i claims, afhængigt af hvordan vi bygger det. Vi skal afklare, hvordan en bruger bliver admin, og hvordan man logger ind som admin i admin-panelet.
    \item Når en bruger ændres i Keycloak, skal tokenet fornyes, før klienten har de nye claims.
    \item Næste skridt i sprint 3 er onboarding af en bruger og at få helt styr på Keycloak-flowet, herunder admin-login i admin-panelet.
    \item Chrome DevTools kan bruges til at se trafikken mellem applikationen og de andre noder.
    \item \textbf{Keycloak:} Keycloak er stort og har mange features. Efterhånden som vi udvider brugen, skal vi dokumentere, hvilke features vi bruger, og hvilke vi har fravalgt. Eksempler, Michel nævnte: hvorfor man kan spærre JWT-tokens, og IP-tjek.
    \item Rapporten skal begrunde valget af Keycloak ud fra de egenskaber, vi bruger, uanset hvorfor vi oprindeligt valgte det.
    \item Keycloak og resten af systemet skal være deployet på nettet. Det er allerede gjort. Michel anbefaler ellers Heroku eller Aarhus Universitets Kubernetes-servere.
    \item Når Keycloak ligger på en offentlig server, skal vi forhindre zombie-brugere.
\end{itemize}

\subsection*{2.1 Diagrammer}

\begin{itemize}
    \item Vi følger Projekt 4-prototypen og de andre eksempler og laver diagrammerne løbende. Så har vi nok.
    \item Michel går ikke meget op i det formelle. Han lægger mere vægt på, at vi får noget til at virke og selv får hænderne i det.
    \item Diagrammerne skal være nogenlunde UML, men detaljer som farven på en pil er ligegyldige.
    \item Dokumentationen er ikke ligegyldig, men den skal være fornuftig og ikke banal.
    \item Vi skal have et overordnet overblik over projektet. Tidsplanen kan rekonstrueres til sidst; den går Michel ikke op i. Det, der ikke kan rekonstrueres, er forståelsen af hele flowet i applikationen og det tekniske.
    \item Vi skal løbende dokumentere, hvordan de forskellige flows virker mellem noderne, fx backend og admin-panelet i React. Det gælder fx flowet, hvor noget gemmes, både for admin-panelet og for kunde-appen.
    \item Flowdiagrammerne behøver ikke alle stå i rapporten, men det skal flettes ind et sted, så det fremgår, at vi har forstået, hvad der foregår i applikationen.
\end{itemize}

\subsection*{2.2 Kilder til kursusmateriale}

\begin{itemize}
    \item Vi skal kun henvise til dokumentation eller litteratur, ikke til slides på Brightspace. Det skal være en original reference.
    \item Eksempel: HTTP-protokollen henvises til dokumentationen for HTTP-protokollen.
    \item Står en vigtig pointe på et slide, finder vi den originale kilde i bogen eller lignende.
    \item Almen softwareviden kræver ingen reference.
\end{itemize}

\section*{Beslutninger}

\begin{itemize}
    \item Rapporten henviser til originale kilder (dokumentation og litteratur) og ikke til slides på Brightspace.
    \item Diagrammerne følger Projekt 4-prototypen og laves løbende i nogenlunde UML.
    \item Sprint 3 fokuserer på onboarding af en bruger og på Keycloak-flowet, herunder claims, roller og admin-login.
    \item Credentials går aldrig gennem backend, som kun håndterer tokens.
\end{itemize}

\section*{Opgaver}

\begin{itemize}
    \item Opgaverne blev ikke fordelt på mødet.
\end{itemize}

\section*{Næste møde}

\begin{itemize}
    \item Onsdag 7/10: sprint review med demo, retrospective og sprint planning for sprint 3.
    \item Næste vejledermøde blev ikke aftalt.
\end{itemize}

\end{document}
